\begin{textsample}{POS Dim 1 – vem\_ed\_28 – Score 22.00 – vem\_ed\_28/t146.txt}  \label{ex:f1_pos_005}
Leitores Aos Osvaldo Succi Jr.
Coordenador PCIs No \textbf{número} \textbf{anterior} deste informativo VEm (janeiro / fevereiro 2025), sintetizamos as apresentações da Jornada de PCIs 2024, realizada em dezembro do ano passado.
Em 20 de fevereiro de 2025, convidamos os professores de Fatecs \textbf{que} apresentaram seus relatos de experiência na Jornada a \textbf{produzirem} textos para a seção “Artigo de Opinião”.
Como resultado, temos as \textbf{contribuições} do Núcleo de Estudos da Linguagem da Fatec (NELF) Bragança Paulista; da professora de inglês Talita Annunciato Rodrigues (Fatec Indaiatuba) e do professor de espanhol Odenildo França Almeida (Fatec Ipiranga).
O texto do NELF Bragança Paulista ressalta \textbf{que} o uso de diversos \textbf{idiomas} (português, inglês, espanhol) nos PCIs / Cesu contribui para a melhoria das \textbf{competências} \textbf{linguísticas} dos \textbf{futuros} tecnólogos, o \textbf{que} \textbf{é} imprescindível para a atuação \textbf{profissional} globalizada.
Talita Rodrigues descreve a experiência de \textbf{aprendizado} intercultural no PCI / Cesu na área de marketing internacional realizado entre a Fatec Indaiatuba e a University of Minnesota Crookston.
Odenildo França Almeida relata o projeto “Escrituras Narrativas”, desenvolvido entre Fatec Ipiranga e Uniminuto desde 2023.
Nas quatro edições, houve PCIs com \textbf{cronogramas} \textbf{variando} entre cinco semanas (“muito corrido”) e oito semanas (“duração \textbf{ideal}”, no entendimento do autor).
Odenildo \textbf{destaca} \textbf{que} a avaliação pós-projeto, feita pelos professores envolvidos, \textbf{é} “\textbf{essencial} para \textbf{identificar} \textbf{pontos} a \textbf{serem} ajustados e resgatar práticas \textbf{bem-sucedidas}”.
Também temos neste \textbf{número}, Selma de Fátima Grossi e Lígia Grandi \textbf{que} escrevem sobre a \textbf{importância} da sinergia entre professores de \textbf{disciplinas} \textbf{profissionais} (\textbf{conteúdos} técnicos) e de \textbf{idiomas} estrangeiros para o sucesso dos Projetos Colaborativos Internacionais (PCIs / Cesu).
Quer contribuir você também?
Acesse https://publicacoescesu.cps.sp.gov.br/VEm/about/submissions Dúvidas?
Escreva para cesu.pci@cps.sp.gov.br.
Boa leitura!

% matched lemmas: anterior, aprendizado, bem-sucedido, competência, conteúdo, contribuição, cronograma, destacar, disciplina, essencial, futuro, ideal, identificar, idioma, importância, linguístico, número, ponto, produzir, profissional, que, ser, variar
\end{textsample}
